% GENERATED FILE. Edit canonical lessons, then run npm run generate:books.
% canonical-source-hash: fnv1a64:7a3fc696190ca75d
% canonical-lessons: GE-C57-gesicht, GE-C57-blut, GE-C57-haut, GE-C57-haar, GE-C57-bauch

\chapter{Face, Blood, Skin, Hair, Belly}
\label{ch:ge-face-blood-skin-hair-belly}
\hlchaptermodality{57}

\begin{chapteropening}
\textbf{By the end of this chapter:} \emph{I can say the face, blood, skin, a hair and the belly.}

Complete the last lesson of chapter 57: I can say the face, blood, skin, a hair and the belly.
\end{chapteropening}

\section[das Gesicht]{das Gesicht — the face}

\label{lesson:GE-C57-gesicht}



\begin{quote}
Before the new one: say the German for a shoulder, then the German for a knee.
\end{quote}

\subsection*{You'll want to know: das Gesicht}
\textbf{das Gesicht} — \textquotedblleft{}the face\textquotedblright{}.

\begin{cousinweb}
\textbf{Gesicht} $\leftarrow$ Middle High German \emph{gesiht}, \textquotedblleft{}sight\textquotedblright{}: \emph{ge-} on the stem of \emph{sehen}, \textquotedblleft{}to see\textquotedblright{}, which EtymDB 2.1 traces to Proto-Germanic \emph{*sehwaną}. The face is what is seen; English \emph{sight} comes from the same verb.
\end{cousinweb}

\begin{grammarlens}[title={What you've built}]
One more word for this chapter.
\end{grammarlens}

\subsection*{Your turn}
\begin{itemize}
\raggedright
  \item \emph{Say it:} \emph{das Gesicht}
  \item \emph{Say it:} \emph{das Gesicht}, once more
  \item \emph{Say it:} say \emph{das Gesicht}
  \item \emph{Recall:} say the German for a shoulder, then the German for a knee, then say \emph{das Gesicht} again
\end{itemize}

\begin{tcolorbox}[breakable,colback=teal!4,colframe=teal!35!black,title={Before you move on}]
What does das Gesicht mean? (\textquotedblleft{}The face\textquotedblright{}.) Say it once more.
\end{tcolorbox}
\section[das Blut]{das Blut — blood}

\label{lesson:GE-C57-blut}



\begin{quote}
Before the new one: say the German for a knee, then the German for the face.
\end{quote}

\subsection*{You'll want to know: das Blut}
\textbf{das Blut} — \textquotedblleft{}blood\textquotedblright{}.

\begin{cousinweb}
\textbf{Blut} $\leftarrow$ Old High German \emph{bluot} $\leftarrow$ Proto-Germanic \emph{*blōþą} — inherited, the same word as English \emph{blood}. (EtymDB 2.1)
\end{cousinweb}

\begin{grammarlens}[title={What you've built}]
One more word for this chapter.
\end{grammarlens}

\subsection*{Your turn}
\begin{itemize}
\raggedright
  \item \emph{Say it:} \emph{das Blut}
  \item \emph{Say it:} \emph{das Blut}, once more
  \item \emph{Say it:} say \emph{das Blut}
  \item \emph{Recall:} say the German for a knee, then the German for the face, then say \emph{das Blut} again
\end{itemize}

\begin{tcolorbox}[breakable,colback=teal!4,colframe=teal!35!black,title={Before you move on}]
What does das Blut mean? (\textquotedblleft{}Blood\textquotedblright{}.) Say it once more.
\end{tcolorbox}
\section[die Haut]{die Haut — skin}

\label{lesson:GE-C57-haut}



\begin{quote}
Before the new one: say the German for the face, then the German for blood.
\end{quote}

\subsection*{You'll want to know: die Haut}
\textbf{die Haut} — \textquotedblleft{}skin\textquotedblright{}.

\begin{cousinweb}
\textbf{Haut} $\leftarrow$ Old High German \emph{hūt} $\leftarrow$ Proto-Germanic \emph{*hūdiz}, \textquotedblleft{}hide\textquotedblright{} — inherited, the same word as English \emph{hide}, an animal's skin. (EtymDB 2.1)
\end{cousinweb}

\begin{grammarlens}[title={What you've built}]
One more word for this chapter.
\end{grammarlens}

\subsection*{Your turn}
\begin{itemize}
\raggedright
  \item \emph{Say it:} \emph{die Haut}
  \item \emph{Say it:} \emph{die Haut}, once more
  \item \emph{Say it:} say \emph{die Haut}
  \item \emph{Recall:} say the German for the face, then the German for blood, then say \emph{die Haut} again
\end{itemize}

\begin{tcolorbox}[breakable,colback=teal!4,colframe=teal!35!black,title={Before you move on}]
What does die Haut mean? (\textquotedblleft{}Skin\textquotedblright{}.) Say it once more.
\end{tcolorbox}
\section[das Haar]{das Haar — a hair; hair}

\label{lesson:GE-C57-haar}



\begin{quote}
Before the new one: say the German for blood, then the German for skin.
\end{quote}

\subsection*{You'll want to know: das Haar}
\textbf{das Haar} — \textquotedblleft{}a hair; hair\textquotedblright{}.

\begin{cousinweb}
\textbf{Haar} $\leftarrow$ Old High German \emph{hār} $\leftarrow$ Proto-Germanic \emph{*hērą} — inherited, the same word as English \emph{hair}. (EtymDB 2.1)
\end{cousinweb}

\begin{grammarlens}[title={What you've built}]
One more word for this chapter.
\end{grammarlens}

\subsection*{Your turn}
\begin{itemize}
\raggedright
  \item \emph{Say it:} \emph{das Haar}
  \item \emph{Say it:} \emph{das Haar}, once more
  \item \emph{Say it:} say \emph{das Haar}
  \item \emph{Recall:} say the German for blood, then the German for skin, then say \emph{das Haar} again
\end{itemize}

\begin{tcolorbox}[breakable,colback=teal!4,colframe=teal!35!black,title={Before you move on}]
What does das Haar mean? (\textquotedblleft{}A hair; hair\textquotedblright{}.) Say it once more.
\end{tcolorbox}
\section[der Bauch]{der Bauch — the belly}

\label{lesson:GE-C57-bauch}



\begin{quote}
Before the new one: say the German for skin, then the German for a hair; hair.
\end{quote}

\subsection*{You'll want to know: der Bauch}
\textbf{der Bauch} — \textquotedblleft{}the belly\textquotedblright{}.

\begin{cousinweb}
\textbf{Bauch} $\leftarrow$ Old High German \emph{būh} $\leftarrow$ Proto-Germanic \emph{*būkaz}, \textquotedblleft{}belly\textquotedblright{} — inherited. Old English \emph{būc}, \textquotedblleft{}belly\textquotedblright{}, is the same word. (EtymDB 2.1)
\end{cousinweb}

\begin{grammarlens}[title={What you've built}]
One more word for this chapter.
\end{grammarlens}

\subsection*{Your turn}
\begin{itemize}
\raggedright
  \item \emph{Say it:} \emph{der Bauch}
  \item \emph{Say it:} \emph{der Bauch}, once more
  \item \emph{Say it:} say \emph{der Bauch}
  \item \emph{Recall:} say the German for skin, then the German for a hair; hair, then say \emph{der Bauch} again
\end{itemize}

\begin{tcolorbox}[breakable,colback=teal!4,colframe=teal!35!black,title={Before you move on}]
What does der Bauch mean? (\textquotedblleft{}The belly\textquotedblright{}.) Say it once more.
\end{tcolorbox}
